\chapter{Introdução à Lógica Proposicional}

\begin{questao}{Questão 1}

Qual das seguintes declarações abaixo \underline{\textbf{não}} é uma proposição?

\begin{alternativas}
    \item A UFJ é uma universidade.
    \item Bolsonaro é o presidente do Brasil.
    \item Qual horário o ônibus passará hoje?
    \item Lula é o presidente do Brasil.
\end{alternativas}

\end{questao}

\begin{comentario_questao}

    \comentarioitem{A}{Esta alternativa é uma proposição, pois é uma declaração de algo que pode ser verdadeiro ou falso.}

    \comentarioitem{B}{Também é uma proposição, pois possui valor lógico bem definido.}

    \comentarioitem{C}{Não é uma proposição, pois trata-se de uma {\bf pergunta}, não podendo ser classificada como verdadeira ou falsa.}

    \comentarioitem{D}{É uma proposição, embora o valor lógico possa ser discutido dependendo do contexto temporal.}

    \respostacorreta{C}

\end{comentario_questao}

\newpage

\begin{questao}{Questão 2}

``Uma mesma pessoa não pode ser, ao mesmo tempo, culpada e inocente em relação ao mesmo fato''.  

Qual o princípio da Lógica que melhor se encaixa no espírito desta declaração?

\begin{alternativas}
    \item Princípio do Terceiro Excluído
    \item Princípio da Casa de Pombos
    \item Princípio da Não-Contradição
    \item Princípio da Não Interferência
\end{alternativas}

\end{questao}

\begin{comentario_questao}

    \comentarioitem{A}{O Princípio do Terceiro Excluído afirma que toda proposição é ou verdadeira ou falsa, não havendo uma terceira possibilidade. Contudo, a afirmação apresentada não trata da exaustividade entre verdade e falsidade, mas sim da impossibilidade de que duas condições opostas ocorram simultaneamente.}

    \comentarioitem{B}{O Princípio da Casa de Pombos é um resultado da combinatória e não possui relação com a análise lógica de proposições contraditórias. Portanto, não se aplica ao contexto da questão.}

    \comentarioitem{C}{O Princípio da Não-Contradição estabelece que uma proposição e sua negação não podem ser verdadeiras ao mesmo tempo e sob as mesmas condições. A frase afirma exatamente isso ao indicar que uma mesma pessoa não pode ser, simultaneamente, culpada e inocente em relação ao mesmo fato.}

    \comentarioitem{D}{O Princípio da Não Interferência não é um princípio clássico da lógica formal, não sendo adequado como resposta para a questão proposta.}

    \respostacorreta{C}

\end{comentario_questao}

\newpage

\begin{questao}{Questão 3}

Qual das seguintes declarações abaixo \underline{\textbf{não}} é uma proposição composta?

\begin{alternativas}
    \item A UFJ é uma universidade e uma instituição.
    \item Lula é o presidente do Brasil até o final do mandato.
    \item Se a leitura do material é feita, então a aula fica mais legal.
    \item Bolsonaro é o presidente do Brasil ou ele é um presidiário.
\end{alternativas}

\end{questao}

\begin{comentario_questao}

    \comentarioitem{A}{Esta é uma proposição composta, pois é formada pela conjunção de duas proposições simples: ``A UFJ é uma universidade'' e ``A UFJ é uma instituição'', unidas pelo conectivo lógico ``e''.}

    \comentarioitem{B}{Esta não é uma proposição composta, pois consiste em uma única afirmação, sem a presença de conectivos lógicos que unam duas ou mais proposições. Trata-se, portanto, de uma proposição simples. ``até o final do mandato'' é um adverbio de tempo que modifica a proposição, mas não a torna composta.}

    \comentarioitem{C}{Esta é uma proposição composta do tipo condicional, formada por duas proposições simples ligadas pelo conectivo ``se... então''.}

    \comentarioitem{D}{Esta é uma proposição composta, pois envolve uma disjunção entre duas proposições simples, conectadas pelo operador lógico ``ou''.}

    \respostacorreta{B}

\end{comentario_questao}

\newpage

\begin{questao}{Questão 4}

As proposições moleculares são formadas por:

\begin{alternativas}
    \item valores lógicos e conectivos
    \item proposições simples e valores lógicos
    \item letras proposicionais e fórmulas
    \item proposições simples e conectivos
    
\end{alternativas}

\end{questao}

\begin{comentario_questao}

    \comentarioitem{A}{Valores lógicos (verdadeiro ou falso) são atribuídos às proposições, mas não são elementos estruturais que formam proposições moleculares. Estas são construídas antes da atribuição de valores lógicos.}

    \comentarioitem{B}{Proposições simples fazem parte da construção de proposições moleculares, mas os valores lógicos não são componentes estruturais da proposição, apenas qualificações atribuídas a ela.}

    \comentarioitem{C}{Letras proposicionais são utilizadas para representar proposições simples, e fórmulas são um outro nome para proposições moleculares. Assim, essa alternativa não descreve corretamente a estrutura básica de formação das proposições moleculares.}

    \comentarioitem{D}{Proposições moleculares são formadas pela combinação de proposições simples por meio de conectivos lógicos, como ``e'', ``ou'', ``se... então'', entre outros. Essa é exatamente a definição de proposição composta.}

    \respostacorreta{D}

\end{comentario_questao}

\newpage

\begin{questao}{Questão 5}

Quantos valores lógicos diferentes possíveis uma proposição composta formada por 4 proposições simples pode ter?

\begin{alternativas}
    \item 2
    \item 8
    \item 16
    \item 32
\end{alternativas}

\end{questao}

\begin{comentario_questao}

    \comentarioitem{A}{Uma proposição composta formada por 4 proposições simples pode ter mais de 2 valores lógicos possíveis, dependendo da combinação dos valores das proposições simples. Seriam 2 valores lógicos possíveis apenas se houvesse 1 proposição simples, pois $2^1 = 2$. Com 4 proposições simples, o número de combinações é maior.}

    \comentarioitem{B}{Com 4 proposições simples, cada uma podendo ser verdadeira ou falsa, o número total de combinações possíveis é $2^4 = 16$. Portanto, a resposta correta não é 8.}

    \comentarioitem{C}{Cada proposição simples pode assumir 2 valores (verdadeiro ou falso), e com 4 proposições simples, o número total de combinações é $2^4 = 16$. Esta é a resposta correta.}

    \comentarioitem{D}{32 valores lógicos possíveis seriam o caso se houvesse 5 proposições simples, pois $2^5 = 32$. Com apenas 4 proposições simples, o número correto é 16.}

    \respostacorreta{C}

\end{comentario_questao}

\newpage

\begin{questao}{Questão 6}

Se $V(p) = F$ e $p$: ``O boné é mais caro do que o chapéu'', então é \underline{\textbf{impossível}} afirmar que:

\begin{alternativas}
    \item o boné custa o mesmo valor do chapéu.
    \item o boné custa R\$ 35,00 e o chapéu custa R\$ 40,00.
    \item o boné custa metade do preço do chapéu.
    \item o boné custa R\$ 15,00 e o chapéu custa R\$ 10,00.
\end{alternativas}

\end{questao}

\begin{comentario_questao}

    \comentarioitem{A}{Se $p$ é falsa, isso significa que a afirmação ``O boné é mais caro do que o chapéu'' não é verdadeira. Portanto, o boné não pode ser mais caro do que o chapéu, mas pode ser do mesmo valor ou mais barato. Assim, afirmar que o boné custa o mesmo valor do chapéu é possível.}

    \comentarioitem{B}{Se $p$ é falsa, isso significa que o boné não é mais caro do que o chapéu. Portanto, é possível que o boné custe R\$ 35,00 e o chapéu custe R\$ 40,00, pois isso não contradiz a falsidade de $p$.}

    \comentarioitem{C}{Se $p$ é falsa, isso significa que o boné não é mais caro do que o chapéu. Portanto, é possível que o boné custe metade do preço do chapéu, pois isso também não contradiz a falsidade de $p$.}

    \comentarioitem{D}{Se $p$ é falsa, isso significa que o boné não é mais caro do que o chapéu. Portanto, afirmar que o boné custa R\$ 15,00 e o chapéu custa R\$ 10,00 seria impossível, pois isso implicaria que o boné é mais barato do que o chapéu, contradizendo a falsidade de $p$.}

    \respostacorreta{D}
\end{comentario_questao}
